%
%Academic CV LaTeX Template
% Author: Dubasi Pavan Kumar
% LinkedIn: https://www.linkedin.com/in/im-pavankumar/
% License: MIT
%
% For errors, suggestions, or improvements, please contact:
% Email: pavankumard.pg19.ma@nitp.ac.in
%



\documentclass[a4paper,11pt]{article}

% Package imports
\usepackage{latexsym}
\usepackage{xcolor}
\usepackage{float}
\usepackage{ragged2e}
\usepackage[empty]{fullpage}
\usepackage{wrapfig}
\usepackage{lipsum}
\usepackage{tabularx}
\usepackage{titlesec}
\usepackage{geometry}
\usepackage{marvosym}
\usepackage{verbatim}
\usepackage{enumitem}
\usepackage{fancyhdr}
\usepackage{multicol}
\usepackage{graphicx}
\usepackage{cfr-lm}
\usepackage[T1]{fontenc}
\usepackage{fontawesome5}

% Color definitions
\definecolor{darkblue}{RGB}{0,0,139}

% Page layout
\setlength{\multicolsep}{0pt} 
\pagestyle{fancy}
\fancyhf{} % clear all header and footer fields
\fancyfoot{}
\renewcommand{\headrulewidth}{0pt}
\renewcommand{\footrulewidth}{0pt}
\geometry{left=1.4cm, top=0.8cm, right=1.2cm, bottom=1cm}
\setlength{\footskip}{5pt} % Addressing fancyhdr warning

% Hyperlink setup (moved after fancyhdr to address warning)
\usepackage[hidelinks]{hyperref}
\hypersetup{
    colorlinks=true,
    linkcolor=darkblue,
    filecolor=darkblue,
    urlcolor=darkblue,
}

% Custom box settings
\usepackage[most]{tcolorbox}
\tcbset{
    frame code={},
    center title,
    left=0pt,
    right=0pt,
    top=0pt,
    bottom=0pt,
    colback=gray!20,
    colframe=white,
    width=\dimexpr\textwidth\relax,
    enlarge left by=-2mm,
    boxsep=4pt,
    arc=0pt,outer arc=0pt,
}

% URL style
\urlstyle{same}

% Text alignment
\raggedright
\setlength{\tabcolsep}{0in}

% Section formatting
\titleformat{\section}{
  \vspace{-4pt}\scshape\raggedright\large
}{}{0em}{}[\color{black}\titlerule \vspace{-7pt}]

% Custom commands
\newcommand{\resumeItem}[2]{
  \item{
    \textbf{#1}{\hspace{0.5mm}#2 \vspace{-0.5mm}}
  }
}

\newcommand{\resumePOR}[3]{
\vspace{0.5mm}\item
    \begin{tabular*}{0.97\textwidth}[t]{l@{\extracolsep{\fill}}r}
        \textbf{#1}\hspace{0.3mm}#2 & \textit{\small{#3}} 
    \end{tabular*}
    \vspace{-2mm}
}

\newcommand{\resumeSubheading}[4]{
\vspace{0.5mm}\item
    \begin{tabular*}{0.98\textwidth}[t]{l@{\extracolsep{\fill}}r}
        \textbf{#1} & \textit{\footnotesize{#4}} \\
        \textit{\footnotesize{#3}} &  \footnotesize{#2}\\
    \end{tabular*}
    \vspace{-2.4mm}
}

\newcommand{\resumeProject}[4]{
\vspace{0.5mm}\item
    \begin{tabular*}{0.98\textwidth}[t]{l@{\extracolsep{\fill}}r}
        \textbf{#1} & \textit{\footnotesize{#3}} \\
        \footnotesize{#2} & \footnotesize{#4}
    \end{tabular*}
    \vspace{-2.4mm}
}

\newcommand{\resumeSubItem}[2]{\resumeItem{#1}{#2}\vspace{-4pt}}

\renewcommand{\labelitemi}{$\vcenter{\hbox{\tiny$\bullet$}}$}
\renewcommand{\labelitemii}{$\vcenter{\hbox{\tiny$\circ$}}$}

\newcommand{\resumeSubHeadingListStart}{\begin{itemize}[leftmargin=*,labelsep=1mm]}
\newcommand{\resumeHeadingSkillStart}{\begin{itemize}[leftmargin=*,itemsep=1.7mm, rightmargin=2ex]}
\newcommand{\resumeItemListStart}{\begin{itemize}[leftmargin=*,labelsep=1mm,itemsep=0.5mm]}

\newcommand{\resumeSubHeadingListEnd}{\end{itemize}\vspace{2mm}}
\newcommand{\resumeHeadingSkillEnd}{\end{itemize}\vspace{-2mm}}
\newcommand{\resumeItemListEnd}{\end{itemize}\vspace{-2mm}}
\newcommand{\cvsection}[1]{%
\vspace{2mm}
\begin{tcolorbox}
    \textbf{\large #1}
\end{tcolorbox}
    \vspace{-4mm}
}

\newcolumntype{L}{>{\raggedright\arraybackslash}X}%
\newcolumntype{R}{>{\raggedleft\arraybackslash}X}%
\newcolumntype{C}{>{\centering\arraybackslash}X}%

% Commands for icon sizing and positioning
\newcommand{\socialicon}[1]{\raisebox{-0.05em}{\resizebox{!}{1em}{#1}}}
\newcommand{\ieeeicon}[1]{\raisebox{-0.3em}{\resizebox{!}{1.3em}{#1}}}

% Font options
\newcommand{\headerfonti}{\fontfamily{phv}\selectfont} % Helvetica-like (similar to Arial/Calibri)
\newcommand{\headerfontii}{\fontfamily{ptm}\selectfont} % Times-like (similar to Times New Roman)
\newcommand{\headerfontiii}{\fontfamily{ppl}\selectfont} % Palatino (elegant serif)
\newcommand{\headerfontiv}{\fontfamily{pbk}\selectfont} % Bookman (readable serif)
\newcommand{\headerfontv}{\fontfamily{pag}\selectfont} % Avant Garde-like (similar to Trebuchet MS)
\newcommand{\headerfontvi}{\fontfamily{cmss}\selectfont} % Computer Modern Sans Serif
\newcommand{\headerfontvii}{\fontfamily{qhv}\selectfont} % Quasi-Helvetica (another Arial/Calibri alternative)
\newcommand{\headerfontviii}{\fontfamily{qpl}\selectfont} % Quasi-Palatino (another elegant serif option)
\newcommand{\headerfontix}{\fontfamily{qtm}\selectfont} % Quasi-Times (another Times New Roman alternative)
\newcommand{\headerfontx}{\fontfamily{bch}\selectfont} % Charter (clean serif font)

\begin{document}
\headerfontiii

% Header
\begin{center}
    {\Huge\textbf{Anbang Wang}}
\end{center}
\vspace{-6mm}

\begin{center}
    \small{
    +86 18853138398 | \href{mailto:awangas@connect.ust.hk}{awangas@connect.ust.hk}
    }
\end{center}
\vspace{-6mm}


\begin{center}
    \small{
    \socialicon{\faLinkedin} \href{https://www.linkedin.com/in/anbang-wang/}{Wang Anbang} | 
    \socialicon{\faGithub} \href{https://github.com/Aawangas}{Aawangas}
    }
\end{center}
\vspace{-6mm}
\begin{center}
    \small{Hong Kong SAR, China}
\end{center}

\vspace{-4mm}

\section{\textbf{Research Interest}}
\vspace{1mm}
\small{
My research interests primarily focus on computer vision and machine learning. Specifically, I am particularly interested in developing fundamental methods to fully leverage large-scale world knowledge. My recent interests are particularly centered on 3D vision.
}
\vspace{-2mm}


\section{\textbf{Education}}
\vspace{-0.4mm}
\resumeSubHeadingListStart

\resumeSubheading
{The Hong Kong University of Science and Technology}{Hong Kong SAR, China}
{BSc in data science and technology}{2023/9 - 2027/9 (expected)}
\resumeItemListStart
\item GPA: 4.035/4.3 (top 2\%)
\resumeItemListEnd

\resumeSubheading
{École Polytechnique Fédérale de Lausanne}{Lausanne, Switzerland}
{Exchange Program in Computer Science section}{2026/2 - 2026/7 (expected)}


\resumeSubHeadingListEnd
\vspace{-4mm}

\section{\textbf{Publications}}

\resumeSubHeadingListStart

\resumeProject
  {SK-Adapter: Skeleton-Based Structural Control for Native 3D Generation}
  {\begin{tabular}[t]{@{}l@{}}\textbf{Anbang Wang}*, Yuzhuo Ao*, Shangzhe Wu, Chi-Keung Tang \\ \textit{* Equal contribution.}\end{tabular}}
  {ECCV 2026}{}

\resumeProject
  {ReasonNavi: Human-Inspired Global Map Reasoning for Zero-Shot Embodied Navigation}
  {\begin{tabular}[t]{@{}l@{}}Yuzhuo Ao*, \textbf{Anbang Wang}*, Yu-Wing Tai, Chi-Keung Tang \\ \textit{* Equal contribution.}\end{tabular}}
  {arXiv}
  {}

\resumeProject
  {MedSapiens: Taking a Pose to Rethink Medical Imaging Landmark Detection}
  {\begin{tabular}[t]{@{}l@{}}Marawan Elbatel*, \textbf{Anbang Wang}*, Keyuan Liu, Kaouther Mouheb, Enrique Almar-Munoz, Lizhuo Lin, Yanqi Yang, \\ Karim Lekadir, Xiaomeng Li \\ \textit{* Equal contribution.}\end{tabular}}
  {MedIA 2026}{}

\resumeProject
  {Geometric-Guided Few-Shot Dental Landmark Detection with Human-Centric Foundation Model}
  {\begin{tabular}[t]{@{}l@{}}\textbf{Anbang Wang}*, Marawan Elbatel*, Keyuan Liu, Lizhuo Lin, Meng Lan, Yanqi Yang, Xiaomeng Li \\ \textit{* Equal contribution.}\end{tabular}}
  {MICCAI 2025}{}
  
\vspace{-4mm}
\resumeSubHeadingListEnd


\section{\textbf{Projects}}
\vspace{-0.4mm}
\resumeSubHeadingListStart

\resumeProject
{SK-Adapter: Skeleton-Based Structural Control for Native 3D Generation}
{\textbf{Advised by Prof. Chi-Keung Tang and Prof. Shangzhe Wu}}
{2025/11 - 2026/03}
{{}[\href{https://sk-adapter.github.io/}{\textcolor{darkblue}{\faGlobe}}]}

\resumeItemListStart
\item Developed SK-Adapter, a lightweight skeleton-based adapter that injects joint-coordinate and topology tokens into a frozen native 3D generation backbone via cross-attention.
\item 
Introduced skeletons as first-class structural control signals for precise 3D asset generation, reducing ambiguity from text- or image-only prompts.
\item 
Contributed to Objaverse-TMS, a 24k text-mesh-skeleton dataset, and demonstrated robust skeleton-conditioned generation and local 3D editing while preserving geometry and texture quality.
\resumeItemListEnd

\resumeProject
{ReasonNavi: Human-Inspired Global Map Reasoning for Zero-Shot Embodied Navigation}
{\textbf{Advised by Prof. Chi-Keung Tang And Prof. Yu-Wing Tai}}
{2025/06 - 2025/09}
{{}[\href{https://reasonnavi.github.io/}{\textcolor{darkblue}{\faGlobe}}]}
  
\resumeItemListStart
\item Developed ReasonNavi, a framework that leverages MLLMs for global reasoning and a deterministic planner for local navigation, creating a novel reason-then-act paradigm.
\item 
Enabled unified, zero-shot navigation across diverse tasks (object, image, text goals), eliminating the need for task-specific fine-tuning or reinforcement learning.
\item 
Demonstrated superior performance over existing methods in efficiency and reliability.
\resumeItemListEnd

\resumeProject
  {Geometric-Guided Few-Shot Dental Landmark Detection with Human-Centric Foundation Model }
  {\textbf{Advised by Prof. Xiaomeng Li}}
  {2024/12 - 2025/03}
  {{}[\href{https://github.com/xmed-lab/GeoSapiens}{\textcolor{darkblue}{\faGithub}}]}
\resumeItemListStart
  \item
  Engineered a novel few-shot learning framework (GeoSapiens) by creatively adapting a human-centric foundation model (Sapiens) for a challenging medical imaging task.
  \item 
Innovated a novel geometric loss function that encodes anatomical structural priors, significantly improving the model's robustness and precision in low-data scenarios.
\item 
Achieved state-of-the-art (SOTA) results, outperforming the previous leading method by 8.18\% in Success Detection Rate (SDR) at the clinically-critical 0.5mm precision threshold.
\resumeItemListEnd

\resumeProject
  {MedSapiens: Taking a Pose to Rethink Medical Imaging Landmark Detection}
  {\textbf{Advised by Prof. Xiaomeng Li}}
  {2024/11 - 2025/11}
  {{}[\href{https://github.com/xmed-lab/MedSapiens}{\textcolor{darkblue}{\faGithub}}]}
\resumeItemListStart
  \item
  Developed MedSapiens by adapting a large-scale human-centric foundation model (Sapiens) for medical anatomical landmark detection, leveraging its inherent spatial pose localization priors across diverse imaging modalities.
  \item 
Harmonized heterogeneous datasets across four distinct anatomical regions (head, hand, chest, and legs) into a unified training pipeline, utilizing LoRA and heatmap-based decoding to capture shared spatial hierarchies across diverse clinical tasks.
\item 
Established new state-of-the-art (SOTA) performance, outperforming generalist models by 5.26\% and specialist models by 21.81\% in Success Detection Rate (SDR), while achieving a 2.69\% improvement in few-shot generalization on novel tasks.
\resumeItemListEnd

\resumeProject
{Unsupervised Domain Adaptation for 3D Bone Reconstruction from 2D X-rays}
{\textbf{Advised by Prof. Xiaomeng Li}} 
{2024/09 - 2024/11} 

\resumeItemListStart
  \item \textbf{Investigated} the challenge of domain shift in 3D bone reconstruction by building a baseline model utilizing the Pixel-Aligned Implicit Function (PiFU) framework.
  \item \textbf{Implemented} a state-of-the-art unsupervised domain adaptation (UDA) algorithm to enhance the model's generalization capabilities across different X-ray datasets.
  \item \textbf{Aimed to improve} reconstruction accuracy and robustness on out-of-domain data without requiring costly new annotations from the target domain.
\resumeItemListEnd

\resumeProject
{Knowledge Distillation on Skin Lesion Segmentation}
{\textbf{Advised by Prof. Albert C. S. Chung}} 
{2024/01 - 2024/09} 
{{}[\href{https://github.com/Aawangas/KD_skin_lesion_segmentation}{\textcolor{darkblue}{\faGithub}}]}

% 然后在下方列出这个职位下的项目和贡献
\resumeItemListStart
    \item \textbf{Developed} a lightweight segmentation model for skin lesions by designing and implementing a knowledge distillation (KD) pipeline.
    \item \textbf{Engineered} a student-teacher framework where knowledge from a large, complex model was effectively transferred to a compact student network.
\resumeItemListEnd

\resumeSubHeadingListEnd

\section{\textbf{Honors and Awards}}
\vspace{-0.4mm}
 \resumeHeadingSkillStart
  \resumeSubItem{HKUST Scholarship for Continuing Undergraduate Students (highest band)}{}
\resumeSubItem{UROP research travel sponsorship}
    {}
  \resumeSubItem{Dean's List for all semesters}
    {}
  
 \resumeHeadingSkillEnd


\section{\textbf{Teaching Experience}}
\vspace{-0.4mm}
\resumeSubHeadingListStart
\resumeProject
  {Teaching Assistant of COMP 2711, Discrete Math and its Applications in Computer Science}
  {HKUST, CSE department}
  {2024/09 - 2024/12}
  {}
  \vspace{1mm}
\resumeItemListStart
  \item \textbf{Assist} weekly tutorial sessions for undergraduate students, clarifying core concepts in logic, set theory, graph theory, and combinatorics. 
  \item \textbf{Provided} regular one-on-one guidance, helping students debug complex proofs and master challenging assignment problems.
\resumeItemListEnd

\resumeSubHeadingListEnd

\vspace{-6mm}

\section{\textbf{Academic Engagement and Service}} %<-- 使用这个更专业的标题
\resumeSubHeadingListStart

  % 条目一：审稿 (服务性质)
  \item \textbf{Reviewer}, IEEE/CVF Conference on Computer Vision and Pattern Recognition (CVPR), 2026

  \item \textbf{Reviewer}, European Conference on Computer Vision (ECCV), 2026

  \item \textbf{Attendee}, International Conference on Medical Image Computing and Computer Assisted Intervention (MICCAI), Daejeon, Republic of Korea, Oct 2025

\resumeSubHeadingListEnd

\vspace{-6mm}

\section{\textbf{Skills}}
\resumeSubHeadingListStart

  \resumeSubItem{Programming Languages: }
    {Python (Proficient), C/C++ (Proficient)}

  \resumeSubItem{Frameworks \& Libraries: }
    {PyTorch, NumPy, Pandas, Scikit-learn, OpenCV} 

  \resumeSubItem{Developer Tools \& OS: }
    {Git, GitHub, Linux, LaTeX} 

  \resumeSubItem{Languages: }
    {Mandarin (Native), English (Fluent - IELTS 7.5/9.0)}

\resumeSubHeadingListEnd

\end{document}
